\documentclass[11pt]{article}
\usepackage{amsmath,amssymb,booktabs,hyperref}
\usepackage[margin=1in]{geometry}
\newcommand{\D}[2]{\lvert D_{#1}^{#2}\rangle}
\title{Explicit anticoherent subspaces from binary polyhedral forms}
\author{}
\date{}
\begin{document}
\maketitle

\paragraph{Scope.}
The constructions below certify lower bounds on the dimension of an
anticoherent (AC) subspace. Their symmetry-sector dimensions are not asserted
to equal the unrestricted maximum. The accompanying files give normalized,
orthonormal Dicke bases and exact verification at every catalogue qubit
number $N=1,\ldots,20,22,24,35,42$.

\paragraph{A whole-subspace symmetry argument.}
Let $G\subset SU(2)$ be a finite group and let $\chi$ be a one-dimensional
unitary character. Consider the subspace on which every collective rotation
$U_g$ acts as $\chi(g)I$. For $u,v$ in this same character sector and any
operator $T$, one has
\[
 \langle u,Tv\rangle
 =\left\langle u,\frac1{|G|}\sum_{g\in G}U_g^*TU_gv\right\rangle.
\]
Thus if the spin-$\ell$ tensor representation contains no $G$-fixed vector,
every tensor of rank $\ell$ compresses to zero on the entire sector, including
its off-diagonal matrix elements. For the quaternion, tetrahedral,
octahedral, and icosahedral rotation groups, the first positive invariant
harmonic degrees are respectively $2,3,4,6$. Any single character sector is
therefore AC of order at least $1,2,3,5$, respectively. This argument does not
permit combining distinct character sectors.

\paragraph{Invariant forms.}
Use the following homogeneous polynomials:
\begin{align*}
 Q&=x^2y^2,& R&=x^4+y^4,& U&=xy,& V&=x^4-y^4,\\
 F&=xy(x^4-y^4),\\
 G&=x^8+14x^4y^4+y^8,\\
 H&=x^{12}-33x^8y^4-33x^4y^8+y^{12},\\
 A&=x^4+2i\sqrt3x^2y^2+y^4,\\
 I&=xy(x^{10}+11x^5y^5-y^{10}),\\
 J&=x^{20}-228x^{15}y^5+494x^{10}y^{10}+228x^5y^{15}+y^{20},\\
 K&=x^{30}+522x^{25}y^5-10005x^{20}y^{10}
       -10005x^{10}y^{20}-522x^5y^{25}+y^{30}.
\end{align*}
Here the bar on $\overline A$ conjugates coefficients. The elementary
relations
\[
 F^2=Q(R^2-4Q^2),\qquad H^2=G^3-108F^4,
 \qquad K^2=J^3+1728I^5
\]
give free-module presentations of the fixed rings. In particular the
quaternion, tetrahedral, octahedral and icosahedral fixed spaces are the
degree-$N$ pieces of
\[
 \mathbb C[Q,R]\{1,F\},\quad
 \mathbb C[F,G]\{1,H\},\quad
 \mathbb C[G,F^2]\{1,FH\},\quad
 \mathbb C[I,J]\{1,K\},
\]
respectively. The notation $\{P_1,P_2\}$ means a free module with the two
displayed generators. Their Hilbert series in the qubit-count variable $z$
are
\[
 \frac{1+z^6}{(1-z^4)^2},\quad
 \frac{1+z^{12}}{(1-z^6)(1-z^8)},\quad
 \frac{1+z^{18}}{(1-z^8)(1-z^{12})},\quad
 \frac{1+z^{30}}{(1-z^{12})(1-z^{20})}.
\]
The quadratic relations must be imposed; simply counting monomials in three
unrestricted generators overcounts the dimension.

\paragraph{Additional character sectors.}
One quaternion nontrivial character sector, one tetrahedral nontrivial
character sector, and the octahedral sign sector are
\[
 \mathbb C[Q,R]\{U,V\},\qquad
 \mathbb C[F,G]\{A,\overline A^{\,2}\},\qquad
 \mathbb C[G,F^2]\{F,H\}.
\]
Their respective Hilbert numerators are $z^2+z^4$, $z^4+z^8$, and $z^6+z^{12}$,
with the same denominators as the corresponding fixed sectors. The identities
\[
 A\overline A=G,\qquad
 A^3=H+6i\sqrt3F^2,\qquad
 \overline A^{\,3}=H-6i\sqrt3F^2
\]
show that $A$ and $\overline A^2$ have the same tetrahedral character.
They are independent over $\mathbb C(F,G)$: otherwise
$\overline A^2/A=G^2/(H+6i\sqrt3F^2)$ would make $H$ rational in $F,G$,
whereas $G^3-108F^4$ is not a square in $\mathbb C(F,G)$.
Similarly $(H/F)^2=(G^3-108(F^2)^2)/F^2$ has a simple pole along $F^2=0$,
so $F,H$ are independent over $\mathbb C(G,F^2)$.
The character Hilbert series then give spanning. For the finite data set,
independence and the character-averaged dimension are additionally checked
directly by exact arithmetic.

\paragraph{Conversion to Dicke vectors.}
For $p(x,y)=\sum_{r=0}^N p_rx^ry^{N-r}$ set
\[
 \langle p,q\rangle_N=\sum_{r=0}^N
 \frac{\overline{p_r}q_r}{\binom Nr}.
\]
Weighted Gram--Schmidt over $\mathbb Q(i\sqrt3)$ produces orthogonal
polynomials $p_1,\ldots,p_k$ with positive rational squared norms $h_i$.
The corresponding orthonormal Dicke basis is
\[
 |u_i\rangle=\sum_{r=0}^N
 \frac{(-1)^rp_{i,r}}{\sqrt{h_i\binom Nr}}\D Nr.
\]
This convention agrees with the homogeneous Majorana polynomial convention
up to an irrelevant common phase.

\paragraph{Exact verification including cross terms.}
The $t$-qubit reduction of $|u_i\rangle\langle u_j|$, in the symmetric Dicke
basis, has $(a,b)$ entry
\[
 \frac{(-1)^{a+b}\sqrt{\binom ta\binom tb}}{\sqrt{h_ih_j}}
 \sum_{r=0}^{N-t}
 \frac{p_{i,a+r}\overline{p_{j,b+r}}\binom{N-t}r}
      {\binom N{a+r}\binom N{b+r}}.
\]
The AC-subspace condition is that this matrix equal
$\delta_{ij}I_{t+1}/(t+1)$. All radicals outside the sum are nonzero and
can be removed. Thus the exact condition checked is
\[
 \sum_{r=0}^{N-t}
 \frac{p_{i,a+r}\overline{p_{j,b+r}}\binom{N-t}r}
      {\binom N{a+r}\binom N{b+r}}
 =\frac{\delta_{ij}\delta_{ab}h_i}{(t+1)\binom ta}.
\]
The program checks all ordered pairs $i,j$ and all entries $a,b$.
It verifies 65 nonzero spaces and 11,896 reduction entries, including the
first order at which each space fails. An independent checker instead uses
the compressed spin moments $J_+^qJ_z^p$ and obtains the same orders.

\paragraph{Compact examples improving catalogue lower bounds.}
The following polynomial spans are particularly useful; normalized
orthonormal Dicke coefficients are supplied in \texttt{bases.json}:
\[
\begin{array}{ccl}
 (N,t,k)& &\text{polynomial span}\\ \hline
 (20,2,2)&&\operatorname{span}\{F^2G,GH\},\\
 (22,3,1)&&\operatorname{span}\{FG^2\},\\
 (24,2,3)&&\operatorname{span}\{G^3,F^4,F^2H\},\\
 (24,3,2)&&\operatorname{span}\{G^3,F^4\},\\
 (42,2,4)&&\operatorname{span}\{F^3G^3,F^7,FG^3H,F^5H\},\\
 (42,3,2)&&\operatorname{span}\{FG^3H,F^5H\}.
\end{array}
\]
Each stated AC order is the exact order of the displayed entire subspace;
each stated dimension is an attained lower bound on the unrestricted
maximum.

\paragraph{References.}
The invariant construction and its guaranteed AC orders follow
Pereira--Paul-Paddock, \href{https://arxiv.org/abs/1610.05841}{Anticoherent
Subspaces}. The Hilbert-series interpretation follows Suter,
\href{https://arxiv.org/abs/math/0503542}{Quantum affine Cartan matrices,
Poincar\'e series of binary polyhedral groups, and reflection
representations}. Explicit icosahedral forms, in a rotated sign convention,
appear in Crass, \href{https://doi.org/10.3390/sym12010177}{Critically-Finite
Dynamics on the Icosahedron}. No maximality claim from the older
first-order AC-subspace corollary is used.
\end{document}
